% ============================================================
% UNIDAD 4: ORGANIZACIONES AERONÁUTICAS Y GESTIÓN DEL MANTENIMIENTO
% ============================================================
\clearpage
\section*{UNIDAD 4: ORGANIZACIONES AERONÁUTICAS}
\subsection*{Gestión, producción, servicios y economía del mantenimiento}

\begin{recuadro}
\textbf{Objetivo de la unidad:} comprender cómo se organizan las empresas del sector aeronáutico,
cómo se planifica y controla el mantenimiento, cómo se abastecen y comercializan partes y
servicios, qué indicadores se usan para gestionar y cuánto pesa el mantenimiento en la economía
de un operador, sin dejar de lado la calidad, el ambiente y la salud laboral.
\end{recuadro}

% ------------------------------------------------------------
\section*{1. LAS ORGANIZACIONES DEL SECTOR}

\begin{center}
\small
\begin{tabular}{|p{0.27\linewidth}|p{0.14\linewidth}|p{0.47\linewidth}|}
\hline
\textbf{Organización} & \textbf{Norma} & \textbf{Qué hace} \\ \hline
Fabricantes de aeronaves, motores, equipos y partes & RAAC 21 & Diseñan y producen bajo
certificado tipo y de producción. \\ \hline
Explotadores aerocomerciales & RAAC 119, 121, 135 & Transportan pasajeros y carga (líneas
aéreas, taxis aéreos). \\ \hline
Explotadores de trabajo aéreo & RAAC 137 & Aviación agrícola, fotografía, inspección. \\ \hline
Aviación general & RAAC 91 & Aeroclubes, propietarios privados, aviación corporativa. \\ \hline
\textbf{Talleres Aeronáuticos de Reparación (TAR)} & \textbf{RAAC 145} & Mantienen, reparan y
modifican aeronaves y componentes para terceros. \\ \hline
Centros de instrucción & RAAC 141, 142, 147 & Forman pilotos y mecánicos. \\ \hline
Proveedores de servicios de navegación aérea & -- & Control de tránsito, radioayudas, radares
(EANA). \\ \hline
Aeropuertos & RAAC 153 & Infraestructura y servicios en tierra. \\ \hline
Distribuidores de partes & -- & Venden y certifican partes y componentes. \\ \hline
\end{tabular}
\end{center}

\subsubsection*{1.1. La estructura de un operador aéreo}
Un explotador aerocomercial tiene, como mínimo, áreas de \textbf{Operaciones} (vuelos y
tripulaciones), \textbf{Mantenimiento e Ingeniería}, \textbf{Calidad}, \textbf{Seguridad
operacional (SMS)}, \textbf{Comercial}, \textbf{Administración y Finanzas} y \textbf{Recursos
Humanos}. Para cada área clave la autoridad exige un \textbf{responsable designado} con la
calificación adecuada.

Dentro de Mantenimiento conviene distinguir dos funciones: la \textbf{gestión de la
aeronavegabilidad continuada} (decidir qué hay que hacer y cuándo: programa, AD, confiabilidad,
registros) y la \textbf{ejecución del mantenimiento} (hacerlo), que puede hacer el propio operador
o contratar a un TAR.

% ------------------------------------------------------------
\section*{2. PRODUCCIÓN Y SERVICIOS: EL MANTENIMIENTO}

\subsubsection*{2.1. Mantenimiento de línea y de base}
\begin{itemize}
  \item \textbf{De línea:} se hace entre vuelos, en la plataforma: inspecciones de tránsito y
        diarias, solución de fallas reportadas por la tripulación, cambio rápido de componentes
        (LRU, \textit{Line Replaceable Unit}). Ahí trabaja mucho el técnico aviónico.
  \item \textbf{De base o pesado:} se hace en hangar, con el avión fuera de servicio varios días
        o semanas: inspecciones mayores, reparaciones estructurales, modificaciones.
\end{itemize}

\subsubsection*{2.2. El programa de mantenimiento}
Las tareas se programan por intervalos de \textbf{horas de vuelo (FH)}, \textbf{ciclos (FC:
despegue y aterrizaje)} o \textbf{calendario}, lo que ocurra primero. En los aviones de línea se
agrupan en \textbf{chequeos}: los más frecuentes y livianos (``A'') y los más espaciados y
profundos (``C'' y, en algunos programas, ``D''). Los programas modernos se elaboran con la
metodología \textbf{MSG-3}, que analiza las consecuencias de cada falla para decidir qué tarea
conviene y con qué frecuencia.

Los componentes se gestionan con tres criterios:
\begin{itemize}
  \item \textbf{Tiempo límite (\textit{hard time}):} se retiran al cumplir su vida, estén como
        estén.
  \item \textbf{A condición (\textit{on condition}):} se inspeccionan o ensayan periódicamente y
        se retiran si no pasan la prueba.
  \item \textbf{Monitoreo de condición:} se usan hasta que fallan, controlando estadísticamente su
        confiabilidad. Es típico de muchos equipos de aviónica, cuya falla es detectable y no
        compromete la seguridad (hay redundancia).
\end{itemize}

\subsubsection*{2.3. Planificación y programación}
La planificación convierte el programa en trabajo concreto: prevé qué aviones vencen qué tareas,
reserva lugar en el hangar, calcula las \textbf{horas-hombre} necesarias, arma el \textbf{paquete
de trabajo} con sus órdenes y tarjetas, y anticipa partes, herramientas y personal licenciado. Una
buena planificación reduce el tiempo en tierra y evita trabajos apurados (y errores).

% ------------------------------------------------------------
\section*{3. ABASTECIMIENTO Y COMERCIALIZACIÓN}

\subsubsection*{3.1. Selección y abastecimiento de partes}
\begin{itemize}
  \item \textbf{Proveedores aprobados:} el taller u operador mantiene una lista de proveedores
        evaluados (calidad, documentación, plazos). Comprar fuera de ella exige una evaluación
        especial.
  \item \textbf{Tipos de partes:} \textbf{rotables} (componentes reparables con número de serie que
        rotan entre avión y taller, como un transceptor o una computadora), \textbf{reparables} y
        \textbf{consumibles} (que se usan una vez: juntas, filtros, tornillería).
  \item \textbf{Intercambio (\textit{exchange}):} se recibe un componente serviciable y se entrega
        el averiado para su reparación, pagando la reparación más una diferencia. Reduce el
        tiempo de avión en tierra.
  \item \textbf{Inspección de recepción:} verificar número de parte, número de serie, estado y
        documentación (8130-3, EASA Form 1) antes de ingresar la parte al pañol.
  \item \textbf{Stock:} definir mínimos y puntos de pedido según consumo y criticidad; las piezas
        que pueden dejar un avión en tierra se priorizan.
\end{itemize}

\subsubsection*{3.2. Comercialización de servicios y equipos}
Un taller vende capacidad técnica: debe tener claro su \textbf{alcance} (lista de capacidades
aprobada por la ANAC), sus \textbf{tiempos de entrega} (TAT, \textit{turn around time}), sus
\textbf{precios} (horas-hombre más materiales), las \textbf{garantías} que ofrece y los
\textbf{contratos} que firma. En aviónica, además, es común la representación y venta de equipos de
fabricantes extranjeros con su instalación y soporte.

% ------------------------------------------------------------
\section*{4. CONTROL DE GESTIÓN}

El \textbf{control de gestión} compara lo planificado con lo real usando \textbf{indicadores}
(KPI), para detectar desvíos y corregirlos a tiempo.

\begin{center}
\small
\begin{tabular}{|p{0.15\linewidth}|p{0.42\linewidth}|p{0.31\linewidth}|}
\hline
\textbf{Área} & \textbf{Indicador} & \textbf{Qué muestra} \\ \hline
Técnica & \textbf{Confiabilidad de despacho}: \% de salidas sin demora técnica & Si el
mantenimiento cumple su objetivo. \\ \hline
Técnica & \textbf{MTBF} (tiempo medio entre fallas) y \textbf{MTBUR} (entre remociones no
programadas) & Confiabilidad de cada componente; detecta equipos ``problemáticos''. \\ \hline
Técnica & Fallas repetitivas, eventos AOG, reportes de pilotos & Calidad del diagnóstico y
de las reparaciones. \\ \hline
Económica & Costo de mantenimiento por hora de vuelo; presupuesto vs.\ real; valor del stock &
Eficiencia y control del gasto. \\ \hline
Comercial & TAT, cumplimiento de plazos, reclamos, facturación & Satisfacción del cliente. \\
\hline
Personal & Horas-hombre productivas, horas extra, capacitación, ausentismo, rotación &
Cargas de trabajo y fatiga (que se conectan con la seguridad). \\ \hline
\end{tabular}
\end{center}

% ------------------------------------------------------------
\section*{5. LA ECONOMÍA DEL OPERADOR Y EL MANTENIMIENTO}

Para un operador, el \textbf{combustible} suele ser el mayor costo, y el \textbf{mantenimiento}
está entre los principales, junto con el personal y el costo de los aviones (compra o
\textit{leasing}). Pero el mantenimiento también influye en los \textbf{ingresos}: un avión en
tierra no vende pasajes.

\begin{recuadro}
\textbf{Caso práctico (valores hipotéticos).} Un operador regional vuela cada avión
\textbf{2.500 horas por año}. Sus costos anuales de mantenimiento por avión son:
\begin{center}
\begin{tabular}{|p{0.55\linewidth}|r|}
\hline
Mantenimiento de línea (personal, materiales) & USD 300.000 \\ \hline
Chequeos de base (prorrateados por año) & USD 400.000 \\ \hline
Reserva para revisiones de motores (USD 250/h $\times$ 2 motores $\times$ 2.500 h) &
USD 1.250.000 \\ \hline
Componentes y reparaciones (USD 150/h $\times$ 2.500 h) & USD 375.000 \\ \hline
\textbf{Total anual por avión} & \textbf{USD 2.325.000} \\ \hline
\textbf{Costo de mantenimiento por hora de vuelo} & \textbf{USD 930} \\ \hline
\end{tabular}
\end{center}
\textbf{Costo de un día AOG:} si el avión hace 6 vuelos por día con 60 pasajeros y un pasaje
medio de USD 100, un día en tierra significa perder unos \textbf{USD 36.000} de ingresos, más la
reubicación de pasajeros y el impacto en la imagen. Por eso conviene invertir en stock de
componentes críticos, en buena planificación y en técnicos bien capacitados.
\end{recuadro}

% ------------------------------------------------------------
\section*{6. CALIDAD, AMBIENTE Y SALUD LABORAL}

\subsubsection*{6.1. Gestión de la calidad}
La RAAC 145 exige un \textbf{sistema de control de calidad} descripto en el manual del taller:
inspección de recepción, inspección durante y después del trabajo, control de herramientas y
calibraciones, auditorías internas, acciones correctivas. Muchos talleres además se certifican en
\textbf{ISO 9001} o \textbf{AS9110}. Calidad y seguridad operacional se complementan: la calidad
verifica que se cumplan los procedimientos; el SMS gestiona los riesgos (Unidad 6).

\subsubsection*{6.2. Protección ambiental}
\begin{itemize}
  \item \textbf{Residuos peligrosos} (Ley 24.051 y normas provinciales): aceites, combustibles,
        fluidos hidráulicos, solventes, pinturas, baterías y residuos electrónicos deben separarse,
        almacenarse rotulados y entregarse a transportistas y operadores habilitados.
  \item \textbf{Derrames:} bateas, materiales absorbentes y procedimientos de respuesta.
  \item \textbf{Ruido y emisiones:} el Anexo 16 de la OACI fija límites para aeronaves y
        motores.
  \item La norma \textbf{ISO 14001} ofrece un sistema de gestión ambiental certificable.
\end{itemize}

\subsubsection*{6.3. Salud y seguridad en el trabajo}
El marco general es la \textbf{Ley 19.587} de Higiene y Seguridad en el Trabajo y la \textbf{Ley
24.557} de Riesgos del Trabajo (ART). Riesgos típicos de un hangar y de un taller de aviónica:
\begin{center}
\small
\begin{tabular}{|p{0.32\linewidth}|p{0.58\linewidth}|}
\hline
\textbf{Riesgo} & \textbf{Prevención} \\ \hline
Caídas en altura (alas, empenaje) & Plataformas con baranda, arnés, escaleras en buen estado. \\
\hline
Ruido de motores y APU & Protección auditiva, zonas delimitadas. \\ \hline
Electricidad (115 V 400 Hz, 28 V DC, fuentes de tierra) & Bloqueo y etiquetado de
interruptores, verificación de ausencia de tensión. \\ \hline
\textbf{Radiofrecuencia}: radar meteorológico y transmisores & No operar el radar en tierra con
personas o combustible cerca del cono de emisión; respetar distancias del manual. \\ \hline
Químicos (fluidos hidráulicos, solventes, combustible) & Guantes, anteojos, ventilación, hojas de
seguridad. \\ \hline
Descarga electrostática (ESD) & Pulsera y superficies antiestáticas al manipular equipos
sensibles. \\ \hline
Objetos extraños (FOD) & Control de herramientas y limpieza de plataforma. \\ \hline
\end{tabular}
\end{center}

% ------------------------------------------------------------
\section*{7. ACTIVIDADES DE LA UNIDAD}
\begin{enumerate}
  \item \textbf{Plan de mantenimiento:} para una flota simulada de tres aviones con sus horas, ciclos
        y vencimientos, armar el plan anual de chequeos y detectar superposiciones.
  \item \textbf{Control de gestión:} con datos de un año de operación (salidas, demoras, remociones
        de componentes), calcular la confiabilidad de despacho y el MTBUR de dos equipos de
        aviónica, y proponer acciones.
  \item \textbf{Economía del operador:} con el caso práctico, calcular cómo cambia el costo por
        hora si el avión vuela 2.000 o 3.000 horas por año, y cuántos días AOG ``paga'' tener un
        transceptor de repuesto en stock.
  \item \textbf{Esquema de gestión de un emprendimiento aeronáutico:} diseñar la organización de un
        pequeño taller de aviónica (producción, previsión de servicios, abastecimiento y
        comercialización de equipos), con su organigrama, sus indicadores y su lista de
        capacidades.
  \item \textbf{Relevamiento de riesgos laborales y ambientales} del taller de la escuela, con
        propuestas de mejora.
\end{enumerate}
